\chapter{Creio na Comunhão dos Santos}
\chaptermark{Creio na Comunhão dos Santos}

Este capítulo aprofunda o terceiro artigo do Símbolo Apostólico, professando a fé na comunhão dos santos. Desdobramos os parágrafos correspondentes do Catecismo da Igreja Católica (CIC 946--962), e o lugar de Maria na Igreja (963--975), revelando a união dos fiéis em Cristo e a maternidade espiritual da Virgem. Para adultos em processo de iniciação cristã pelo RICA, essa comunhão fortalece a vida sacramental, unindo catecúmenos aos santos na Vigília Pascal.

\section{A Comunhão dos Santos (\S 946--962)}

Após confessar ``a santa Igreja católica'', o Credo adiciona ``a comunhão dos santos'', explicação ulterior: a Igreja é assembleia de todos os santos.

Todos os fiéis formam um corpo: o bem de cada um comunica-se aos outros, tesouro comum presidido por Cristo Cabeça, via sacramentos.

``Comunhão dos santos'' tem duplo sentido: comunhão nas coisas santas (sancta) e entre pessoas santas (sancti); os fiéis alimentam-se do Corpo e Sangue de Cristo para crescer na koinonia do Espírito.

Comunhão na fé: tesouro apostólico partilhado, enriquecido no encontro eclesial.

Comunhão dos sacramentos: frutos para todos, especialmente Eucaristia que une a Deus e aos irmãos; Batismo como porta.

Para adultos a partir dos 20 anos no RICA, a comunhão dos santos no rito de eleição inspira invocação dos santos padroeiros, integrando catecúmenos à rede espiritual da Igreja militante.

\section{Maria, Mãe da Igreja (\S 963--975)}

Maria ocupa lugar singular no mistério da Igreja: Mãe de Deus e do Redentor, mãe dos membros de Cristo pela caridade que gera fiéis.

Sua união com o Filho na salvação manifesta-se da concepção à Cruz: perseverou na fé, unindo-se ao sacrifício, dada como mãe ao discípulo.

Após Ascensão, aids os primórdios da Igreja pela oração, implorando o Espírito em Pentecostes; Assunção singular participa da Ressurreição.

Modelo de fé e caridade para a Igreja, membro preeminente, typus da Igreja exemplar.

Mãe na ordem da graça: coopera na restauração da vida sobrenatural; continua como Advogada, Auxiliadora, Benfeitora, Mediatrix, sem obscurecer mediação de Cristo.

Na catequese para adultos não iniciados, Maria no tempo quaresmal do catecumenato, com rosário e consagração, desperta fiat filial, preparando para a Eucaristia mariana na iniciação pascal.

\section{Conclusão}

Professar a comunhão dos santos e Maria une vivos e celestiais em Cristo. Adultos no RICA, vivam essa koinonia sacramental, com Maria rumo à santidade eterna.